Oui, et petite correction importante :
[
\mathbb E(X_1+\cdots+X_n)=\mathbb E(X_1)+\cdots+\mathbb E(X_n)
]
ne demande **pas** l’indépendance. C’est la **linéarité de l’espérance**.
L’indépendance est nécessaire pour écrire :
[
\mathrm{Var}(X_1+\cdots+X_n)=\mathrm{Var}(X_1)+\cdots+\mathrm{Var}(X_n).
]

Voici une version LaTeX plus rigoureuse et plus contextualisée.

```latex
\documentclass[11pt,a4paper]{article}

%=========================================================
% PACKAGES
%=========================================================
\usepackage[utf8]{inputenc}
\usepackage[T1]{fontenc}
\usepackage[french]{babel}
\usepackage{lmodern}
\usepackage{geometry}
\geometry{margin=2.1cm}
\usepackage{amsmath,amssymb,amsfonts}
\usepackage{mathtools}
\usepackage{siunitx}
\usepackage{xcolor}
\usepackage{tcolorbox}
\usepackage{enumitem}
\usepackage{hyperref}
\usepackage{fancyhdr}
\usepackage{titlesec}
\usepackage{array}
\usepackage{booktabs}
\usepackage{multicol}

%=========================================================
% STYLE
%=========================================================
\hypersetup{
    colorlinks=true,
    linkcolor=blue!60!black,
    urlcolor=blue!60!black
}

\pagestyle{fancy}
\fancyhf{}
\lhead{\textsc{Grand Oral -- Probabilités}}
\rhead{\textsc{Jeux d'argent}}
\cfoot{\thepage}

\titleformat{\section}{\Large\bfseries\color{blue!55!black}}{}{0em}{}
\titleformat{\subsection}{\large\bfseries\color{blue!40!black}}{}{0em}{}
\titleformat{\subsubsection}{\bfseries\color{blue!30!black}}{}{0em}{}

\tcbset{
    arc=2mm,
    boxrule=0.7pt
}

\newtcolorbox{idee}[1]{
    title=\textbf{Idée générale : #1},
    colback=green!5,
    colframe=green!45!black
}

\newtcolorbox{contexte}[1]{
    title=\textbf{Contexte concret : #1},
    colback=orange!5,
    colframe=orange!70!black
}

\newtcolorbox{definitionbox}[1]{
    title=\textbf{Définition : #1},
    colback=blue!4,
    colframe=blue!55!black
}

\newtcolorbox{modele}[1]{
    title=\textbf{Modélisation : #1},
    colback=cyan!4,
    colframe=cyan!55!black
}

\newtcolorbox{equationbox}[1]{
    title=\textbf{Mise en équation : #1},
    colback=yellow!7,
    colframe=orange!80!black
}

\newtcolorbox{preuve}[1]{
    title=\textbf{Justification mathématique : #1},
    colback=purple!4,
    colframe=purple!55!black
}

\newtcolorbox{resolution}[1]{
    title=\textbf{Résolution : #1},
    colback=gray!7,
    colframe=black!60
}

\newtcolorbox{conclusionbox}[1]{
    title=\textbf{Conclusion sur le jeu : #1},
    colback=red!4,
    colframe=red!60!black
}

\newtcolorbox{oral}[1]{
    title=\textbf{Rédaction possible à l'oral : #1},
    colback=white,
    colframe=black!70
}

\newtcolorbox{prepa}[1]{
    title=\textbf{Ouverture vers la prépa : #1},
    colback=teal!4,
    colframe=teal!55!black
}

\newtcolorbox{attention}[1]{
    title=\textbf{Attention : #1},
    colback=red!7,
    colframe=red!80!black
}

%=========================================================
% COMMANDES
%=========================================================
\newcommand{\E}{\mathbb{E}}
\newcommand{\V}{\mathrm{Var}}
\newcommand{\Prob}{\mathbb{P}}
\newcommand{\N}{\mathbb{N}}
\newcommand{\R}{\mathbb{R}}

%=========================================================
% DOCUMENT
%=========================================================
\begin{document}

\begin{center}
    {\Huge \textbf{Probabilités et jeux d'argent}}\\[0.2cm]
    {\Large Peut-on battre le hasard ?}\\[0.2cm]
    {\large Dossier rigoureux pour Grand Oral}\\
    {\large Niveau Terminale approfondi, avec ouvertures vers la prépa}\\[0.4cm]
    \rule{0.85\linewidth}{0.5pt}
\end{center}

\vspace{0.4cm}

\tableofcontents

\newpage

%=========================================================
\section{Introduction générale}

Les jeux d'argent attirent parce qu'ils donnent l'impression qu'une stratégie, une intuition ou une suite de décisions peut permettre de battre le hasard. Un joueur peut penser :

\begin{itemize}
    \item qu'il suffit d'attendre la bonne série ;
    \item qu'après plusieurs pertes, une victoire est plus probable ;
    \item qu'en changeant de jeu, il peut casser une mauvaise dynamique ;
    \item qu'une martingale permet de gagner presque sûrement ;
    \item qu'en jouant tous les jours, il finira forcément gagnant ;
    \item qu'au blackjack, une bonne stratégie peut renverser l'avantage du casino.
\end{itemize}

Le but de ce dossier est de montrer que ces idées peuvent être étudiées rigoureusement avec les probabilités.

\begin{center}
\Large
\textbf{Problématique possible :}\\[0.2cm]
\textit{Les mathématiques permettent-elles de gagner aux jeux d'argent, ou seulement de comprendre pourquoi on perd ?}
\end{center}

\begin{attention}{objectif du document}
Ce document ne donne pas de méthode pour encourager les jeux d'argent. Au contraire, il montre que les jeux d'argent sont souvent défavorables au joueur et que les stratégies populaires reposent fréquemment sur une mauvaise compréhension des probabilités.
\end{attention}

%=========================================================
\section{Préliminaires mathématiques rigoureux}

Avant de traiter les pistes, on fixe les outils mathématiques utilisés.

%---------------------------------------------------------
\subsection{Variable aléatoire}

\begin{definitionbox}{variable aléatoire}
Une variable aléatoire est une fonction qui associe un nombre réel à chaque issue possible d'une expérience aléatoire.

Dans un jeu d'argent, on utilise souvent une variable aléatoire $X$ pour représenter le gain algébrique du joueur.

Le gain algébrique signifie :

\[
X = \text{argent reçu} - \text{mise payée}.
\]

Ainsi, si un joueur mise $1$ euro et reçoit $10$ euros, son gain algébrique est :

\[
X=10-1=9.
\]

S'il ne reçoit rien, son gain algébrique est :

\[
X=0-1=-1.
\]
\end{definitionbox}

%---------------------------------------------------------
\subsection{Espérance}

\begin{definitionbox}{espérance d'une variable aléatoire discrète}
Si une variable aléatoire $X$ prend les valeurs $x_1,x_2,\dots,x_r$ avec les probabilités associées $p_1,p_2,\dots,p_r$, alors son espérance est :

\[
\E(X)=\sum_{i=1}^{r}x_i p_i.
\]

L'espérance représente la moyenne théorique des gains si l'on répétait le jeu un très grand nombre de fois.
\end{definitionbox}

\begin{exemple}
Si :

\[
X=
\begin{cases}
9 & \text{avec probabilité } 0.1,\\
-1 & \text{avec probabilité }0.9,
\end{cases}
\]

alors :

\[
\E(X)=9\times 0.1+(-1)\times 0.9=0.9-0.9=0.
\]

Le jeu est équitable.
\end{exemple}

%---------------------------------------------------------
\subsection{Linéarité de l'espérance}

\begin{preuve}{linéarité de l'espérance}
La propriété fondamentale est :

\[
\E(aX+bY)=a\E(X)+b\E(Y).
\]

En particulier :

\[
\E(X+Y)=\E(X)+\E(Y).
\]

Cette propriété est vraie même si $X$ et $Y$ ne sont pas indépendantes.

Pour comprendre pourquoi, supposons que $X$ et $Y$ soient deux variables aléatoires définies sur un univers fini $\Omega$.

Alors :

\[
\E(X+Y)=\sum_{\omega\in\Omega}(X(\omega)+Y(\omega))\Prob(\{\omega\}).
\]

On distribue la somme :

\[
\E(X+Y)=\sum_{\omega\in\Omega}X(\omega)\Prob(\{\omega\})
+
\sum_{\omega\in\Omega}Y(\omega)\Prob(\{\omega\}).
\]

Donc :

\[
\E(X+Y)=\E(X)+\E(Y).
\]

Par récurrence :

\[
\E(X_1+\cdots+X_n)=\E(X_1)+\cdots+\E(X_n).
\]

Si les $X_i$ ont toutes la même espérance $\mu$, alors :

\[
\E(X_1+\cdots+X_n)=n\mu.
\]
\end{preuve}

\begin{attention}{indépendance inutile pour l'espérance}
Il ne faut pas dire que :

\[
\E(S_n)=n\E(X)
\]

vient de l'indépendance. Cette relation vient de la linéarité de l'espérance. L'indépendance est nécessaire plus tard pour simplifier la variance.
\end{attention}

%---------------------------------------------------------
\subsection{Variance}

\begin{definitionbox}{variance}
La variance mesure la dispersion d'une variable aléatoire autour de son espérance.

Elle est définie par :

\[
\V(X)=\E\left((X-\E(X))^2\right).
\]

On utilise souvent la formule de König-Huygens :

\[
\V(X)=\E(X^2)-\E(X)^2.
\]

L'écart-type est :

\[
\sigma(X)=\sqrt{\V(X)}.
\]
\end{definitionbox}

\begin{preuve}{formule de König-Huygens}
On part de la définition :

\[
\V(X)=\E\left((X-\E(X))^2\right).
\]

On développe :

\[
(X-\E(X))^2=X^2-2X\E(X)+\E(X)^2.
\]

Donc :

\[
\V(X)=\E(X^2-2X\E(X)+\E(X)^2).
\]

Par linéarité :

\[
\V(X)=\E(X^2)-2\E(X)\E(X)+\E(X)^2.
\]

Ainsi :

\[
\V(X)=\E(X^2)-\E(X)^2.
\]
\end{preuve}

%---------------------------------------------------------
\subsection{Indépendance}

\begin{definitionbox}{indépendance de deux événements}
Deux événements $A$ et $B$ sont indépendants si :

\[
\Prob(A\cap B)=\Prob(A)\Prob(B).
\]

Cela signifie que la réalisation de $A$ ne modifie pas la probabilité de réalisation de $B$.
\end{definitionbox}

\begin{definitionbox}{indépendance de variables aléatoires}
Deux variables aléatoires $X$ et $Y$ sont indépendantes si les événements associés à leurs valeurs sont indépendants.

De manière simplifiée, cela signifie que connaître la valeur de $X$ ne donne aucune information sur la valeur de $Y$.
\end{definitionbox}

\begin{exemple}
Deux lancers successifs d'une pièce équilibrée sont indépendants. Si le premier lancer donne pile, la probabilité d'obtenir pile au second lancer reste :

\[
\frac12.
\]
\end{exemple}

%---------------------------------------------------------
\subsection{Variance d'une somme}

\begin{preuve}{variance d'une somme de variables indépendantes}
Pour deux variables aléatoires $X$ et $Y$, on a :

\[
\V(X+Y)=\E((X+Y)^2)-\E(X+Y)^2.
\]

Or :

\[
(X+Y)^2=X^2+2XY+Y^2.
\]

Donc :

\[
\E((X+Y)^2)=\E(X^2)+2\E(XY)+\E(Y^2).
\]

De plus :

\[
\E(X+Y)^2=(\E(X)+\E(Y))^2.
\]

Donc :

\[
\E(X+Y)^2=\E(X)^2+2\E(X)\E(Y)+\E(Y)^2.
\]

Ainsi :

\[
\V(X+Y)
=
\E(X^2)-\E(X)^2
+
\E(Y^2)-\E(Y)^2
+
2(\E(XY)-\E(X)\E(Y)).
\]

On obtient :

\[
\V(X+Y)=\V(X)+\V(Y)+2\mathrm{Cov}(X,Y),
\]

où :

\[
\mathrm{Cov}(X,Y)=\E(XY)-\E(X)\E(Y).
\]

Si $X$ et $Y$ sont indépendantes, alors :

\[
\E(XY)=\E(X)\E(Y).
\]

Donc :

\[
\mathrm{Cov}(X,Y)=0.
\]

Ainsi :

\[
\V(X+Y)=\V(X)+\V(Y).
\]

Par récurrence, si $X_1,\dots,X_n$ sont indépendantes :

\[
\V(X_1+\cdots+X_n)=\V(X_1)+\cdots+\V(X_n).
\]
\end{preuve}

%---------------------------------------------------------
\subsection{Loi binomiale}

\begin{definitionbox}{loi binomiale}
On dit que $X$ suit une loi binomiale de paramètres $n$ et $p$, notée :

\[
X\sim \mathcal{B}(n,p),
\]

si $X$ compte le nombre de succès dans $n$ expériences indépendantes, chacune ayant la même probabilité de succès $p$.

Alors :

\[
\Prob(X=k)=\binom{n}{k}p^k(1-p)^{n-k}.
\]

Son espérance est :

\[
\E(X)=np.
\]

Sa variance est :

\[
\V(X)=np(1-p).
\]
\end{definitionbox}

%---------------------------------------------------------
\subsection{Espérance conditionnelle simplifiée}

\begin{definitionbox}{espérance conditionnelle, version accessible}
Dans certains jeux, comme le blackjack, on connaît une information partielle avant de décider.

Par exemple, on connaît son total de points et la carte visible du croupier.

L'espérance conditionnelle d'une décision est l'espérance du gain en tenant compte de cette information.

On peut l'interpréter comme la moyenne théorique du gain si l'on se trouve dans une situation donnée.
\end{definitionbox}

%---------------------------------------------------------
\subsection{Convergence en probabilité et loi des grands nombres}

\begin{definitionbox}{convergence en probabilité, idée intuitive}
Dire qu'une suite de variables aléatoires $Y_n$ converge en probabilité vers un nombre $\ell$ signifie que $Y_n$ a de très grandes chances d'être proche de $\ell$ lorsque $n$ est grand.

On note :

\[
Y_n \xrightarrow[n\to+\infty]{\mathbb{P}} \ell.
\]
\end{definitionbox}

\begin{prepa}{loi faible des grands nombres}
Si $X_1,X_2,\dots$ sont des variables aléatoires indépendantes, de même loi, d'espérance $\mu$, alors :

\[
\frac{X_1+\cdots+X_n}{n}
\xrightarrow[n\to+\infty]{\mathbb{P}}
\mu.
\]

Cela signifie que la moyenne observée se rapproche de l'espérance théorique quand le nombre de parties devient grand.
\end{prepa}

\newpage

%=========================================================
\section{Piste 1 -- Espérance d'un jeu : pourquoi un gros gain ne suffit pas}

\subsection{Contexte concret}

\begin{contexte}{ticket de loterie}
Un joueur achète un ticket de loterie à $2$ euros. Le ticket annonce : ``Gagnez jusqu'à $100$ euros !''.

Le joueur est attiré par le gros gain possible, mais il oublie une question essentielle : quelle est la probabilité de ce gain ?

Les probabilités permettent de répondre à la vraie question :

\[
\text{Le jeu est-il rentable en moyenne ?}
\]
\end{contexte}

\subsection{Idée générale}

\begin{idee}{espérance}
On modélise le gain du joueur par une variable aléatoire. L'espérance permet de savoir combien le joueur gagne ou perd en moyenne par partie.
\end{idee}

\subsection{Modélisation}

On suppose que le ticket coûte $2$ euros.

Le joueur :
\begin{itemize}
    \item gagne $100$ euros avec probabilité $0.01$ ;
    \item gagne $0$ euro avec probabilité $0.99$.
\end{itemize}

On note $X$ le gain algébrique.

Si le joueur gagne :

\[
X=100-2=98.
\]

S'il perd :

\[
X=0-2=-2.
\]

Donc :

\[
X=
\begin{cases}
98 & \text{avec probabilité }0.01,\\
-2 & \text{avec probabilité }0.99.
\end{cases}
\]

\subsection{Mise en équation}

\[
\E(X)=98\times 0.01+(-2)\times 0.99.
\]

\subsection{Résolution}

\[
\E(X)=0.98-1.98=-1.
\]

Donc :

\[
\E(X)=-1.
\]

\subsection{Conclusion}

\begin{conclusionbox}{ticket de loterie}
Le joueur perd en moyenne $1$ euro par ticket.

Même si le ticket permet parfois de gagner $100$ euros, le jeu est défavorable car la probabilité de ce gain est trop faible.

La formule montre que le vrai critère n'est pas le gain maximal, mais l'espérance.
\end{conclusionbox}

\begin{oral}{rédaction}
Pour étudier ce ticket de loterie, je définis une variable aléatoire $X$ égale au gain algébrique du joueur. Si le joueur gagne, son gain net est $98$ euros, car il faut retirer le prix du ticket. S'il perd, son gain est $-2$ euros. Le calcul de l'espérance donne $-1$. Cela signifie que, sur un grand nombre de tickets, le joueur perd en moyenne $1$ euro par ticket. Le jeu peut donc faire gagner ponctuellement, mais il est perdant en moyenne.
\end{oral}

%=========================================================
\section{Piste 2 -- Répéter un jeu : suite de gains et perte moyenne}

\subsection{Contexte concret}

\begin{contexte}{machine à sous}
Un joueur ne joue pas une seule fois à une machine à sous. Il joue souvent une série de parties : $10$, $100$, parfois davantage.

La question devient alors :

\[
\text{Que devient le gain total après }n\text{ parties ?}
\]
\end{contexte}

\subsection{Idée générale}

\begin{idee}{somme de variables aléatoires}
Chaque partie produit un gain aléatoire. Le gain total est la somme des gains de toutes les parties.
\end{idee}

\subsection{Modélisation}

On note :

\[
X_i
\]

le gain algébrique à la $i$-ème partie.

Le gain total après $n$ parties est :

\[
S_n=X_1+X_2+\cdots+X_n.
\]

On suppose que toutes les parties ont la même espérance :

\[
\E(X_i)=\mu.
\]

\subsection{Mise en équation}

Par linéarité de l'espérance :

\[
\E(S_n)=\E(X_1+\cdots+X_n).
\]

Donc :

\[
\E(S_n)=\E(X_1)+\cdots+\E(X_n).
\]

Comme :

\[
\E(X_i)=\mu,
\]

on obtient :

\[
\E(S_n)=n\mu.
\]

\begin{attention}{rôle de l'indépendance}
Cette égalité ne demande pas l'indépendance des parties. Elle vient uniquement de la linéarité de l'espérance.

L'indépendance sera nécessaire pour calculer proprement la variance de $S_n$.
\end{attention}

\subsection{Résolution dans un cas défavorable}

Supposons que la machine ait une espérance :

\[
\mu=-0.05.
\]

Cela signifie que le joueur perd en moyenne $5$ centimes par partie.

Alors :

\[
\E(S_n)=-0.05n.
\]

Pour $n=100$ :

\[
\E(S_{100})=-5.
\]

Pour $n=1000$ :

\[
\E(S_{1000})=-50.
\]

\subsection{Conclusion}

\begin{conclusionbox}{répétition}
Répéter un jeu défavorable ne le rend pas favorable.

Au contraire :

\[
\E(S_n)=n\mu
\]

montre que la perte moyenne croît linéairement avec le nombre de parties lorsque $\mu<0$.
\end{conclusionbox}

\begin{oral}{rédaction}
Si le joueur répète le même jeu, je note $X_i$ le gain de la $i$-ème partie et $S_n$ le gain total après $n$ parties. Grâce à la linéarité de l'espérance, on obtient $\E(S_n)=n\mu$, où $\mu$ est l'espérance d'une partie. Si cette espérance est négative, alors plus le joueur joue, plus sa perte moyenne augmente. La répétition n'efface pas l'avantage du casino : elle l'amplifie.
\end{oral}

%=========================================================
\section{Piste 3 -- Variance : pourquoi peut-on gagner à court terme malgré un jeu perdant ?}

\subsection{Contexte concret}

\begin{contexte}{illusion des séries gagnantes}
Un joueur peut sortir du casino gagnant après quelques parties, même si le jeu est mathématiquement défavorable.

Cela ne contredit pas l'espérance négative : c'est dû aux fluctuations aléatoires, mesurées par la variance.
\end{contexte}

\subsection{Idée générale}

\begin{idee}{espérance contre dispersion}
L'espérance donne la tendance moyenne. La variance mesure l'ampleur des écarts possibles autour de cette tendance.
\end{idee}

\subsection{Modélisation}

On reprend :

\[
S_n=X_1+\cdots+X_n.
\]

On suppose que les $X_i$ sont indépendantes, de même loi, avec :

\[
\E(X_i)=\mu,
\qquad
\V(X_i)=\sigma^2.
\]

\subsection{Mise en équation}

On a :

\[
\E(S_n)=n\mu.
\]

Comme les variables sont indépendantes :

\[
\V(S_n)=\V(X_1+\cdots+X_n).
\]

Donc :

\[
\V(S_n)=\V(X_1)+\cdots+\V(X_n).
\]

Ainsi :

\[
\V(S_n)=n\sigma^2.
\]

L'écart-type de $S_n$ vaut :

\[
\sigma(S_n)=\sqrt{\V(S_n)}=\sigma\sqrt{n}.
\]

\subsection{Comparaison des ordres de grandeur}

La tendance moyenne est :

\[
n\mu.
\]

Elle est proportionnelle à $n$.

L'écart-type est :

\[
\sigma\sqrt{n}.
\]

Il est proportionnel à $\sqrt{n}$.

Or :

\[
\frac{\sqrt{n}}{n}=\frac{1}{\sqrt{n}}\to 0.
\]

Donc, à long terme, la tendance moyenne domine les fluctuations relatives.

\subsection{Conclusion}

\begin{conclusionbox}{court terme et long terme}
À court terme, un joueur peut gagner grâce aux fluctuations.

Mais à long terme, si l'espérance est négative, la perte moyenne finit par dominer.

La variance explique pourquoi le jeu peut être psychologiquement piégeux : elle entretient l'espoir par des gains ponctuels.
\end{conclusionbox}

\begin{oral}{rédaction}
La variance permet de comprendre pourquoi un jeu défavorable peut quand même donner des gagnants ponctuels. Pour $n$ parties indépendantes, l'espérance totale est proportionnelle à $n$, tandis que l'écart-type est proportionnel à $\sqrt{n}$. Les fluctuations existent donc, mais elles deviennent relativement moins importantes quand $n$ augmente. À long terme, c'est l'espérance qui domine.
\end{oral}

%=========================================================
\section{Piste 4 -- Changer de jeu : stratégie ou illusion ?}

\subsection{Contexte concret}

\begin{contexte}{changer de machine}
Un joueur perd sur une machine à sous. Il se dit : ``Cette machine ne paie pas, je vais en essayer une autre.''

La question est :

\[
\text{Changer de jeu peut-il transformer une situation défavorable en situation favorable ?}
\]
\end{contexte}

\subsection{Idée générale}

\begin{idee}{moyenne pondérée d'espérances}
Si le joueur alterne entre plusieurs jeux, son espérance globale est une moyenne pondérée des espérances de ces jeux.
\end{idee}

\subsection{Modélisation}

On suppose que le joueur peut choisir entre trois jeux :

\[
J_1,\quad J_2,\quad J_3.
\]

Les espérances des gains sont :

\[
\mu_1,\quad \mu_2,\quad \mu_3.
\]

Le joueur choisit :
\begin{itemize}
    \item $J_1$ avec probabilité $q_1$ ;
    \item $J_2$ avec probabilité $q_2$ ;
    \item $J_3$ avec probabilité $q_3$.
\end{itemize}

Avec :

\[
q_1+q_2+q_3=1.
\]

\subsection{Mise en équation}

L'espérance du gain d'une partie est :

\[
\E(Y)=q_1\mu_1+q_2\mu_2+q_3\mu_3.
\]

Plus généralement :

\[
\E(Y)=\sum_{j=1}^{k}q_j\mu_j.
\]

\subsection{Résolution}

Si tous les jeux sont défavorables :

\[
\mu_j<0
\]

pour tout $j$.

Comme :

\[
q_j\geq 0,
\]

alors :

\[
q_j\mu_j\leq 0.
\]

Donc :

\[
\E(Y)<0
\]

dès que le joueur choisit avec probabilité non nulle au moins un jeu strictement défavorable.

\subsection{Conclusion}

\begin{conclusionbox}{changer de jeu}
Changer entre plusieurs jeux défavorables ne crée pas un jeu favorable.

Cela peut modifier le risque, la variance, ou les sensations du joueur, mais pas rendre l'espérance positive.
\end{conclusionbox}

\begin{oral}{rédaction}
On peut modéliser le changement de jeu par un choix aléatoire entre plusieurs jeux. L'espérance globale est alors une moyenne pondérée des espérances des différents jeux. Si toutes ces espérances sont négatives, leur moyenne pondérée reste négative. Changer de jeu ne permet donc pas, à lui seul, de battre l'avantage mathématique du casino.
\end{oral}

%=========================================================
\section{Piste 5 -- Martingale : doubler après chaque perte}

\subsection{Contexte concret}

\begin{contexte}{roulette ou pile-face}
Un joueur pense avoir trouvé une méthode infaillible : miser $1$, puis doubler après chaque perte. Dès qu'il gagne, il récupère toutes ses pertes et gagne $1$ unité.

Cette stratégie est appelée martingale classique.
\end{contexte}

\subsection{Idée générale}

\begin{idee}{petits gains fréquents, grosse perte rare}
La martingale donne beaucoup de petites victoires, mais expose à une perte rare et énorme.
\end{idee}

\subsection{Modélisation}

On considère un jeu où le joueur :
\begin{itemize}
    \item gagne sa mise avec probabilité $p$ ;
    \item perd sa mise avec probabilité $q=1-p$.
\end{itemize}

Pour commencer, on prend le cas équilibré :

\[
p=q=\frac12.
\]

La stratégie est :
\[
1,\;2,\;4,\;8,\;\dots,\;2^r.
\]

\subsection{Mise en équation}

Après $r$ pertes consécutives, la somme déjà perdue est :

\[
1+2+4+\cdots+2^{r-1}.
\]

C'est une somme géométrique :

\[
1+2+\cdots+2^{r-1}=2^r-1.
\]

Si le joueur mise ensuite $2^r$ et gagne, son bilan est :

\[
2^r-(2^r-1)=1.
\]

Donc une victoire après une série de pertes donne toujours un gain net de $1$.

\subsection{Capital nécessaire}

Pour supporter $r$ pertes consécutives, il faut pouvoir payer :

\[
2^r-1.
\]

Si le capital disponible est $C$, il faut :

\[
2^r-1\leq C.
\]

Donc :

\[
2^r\leq C+1.
\]

Ainsi :

\[
r\leq \log_2(C+1).
\]

Le nombre de pertes supportables augmente seulement comme un logarithme du capital, alors que les mises augmentent exponentiellement.

\subsection{Probabilité de ruine}

Si le joueur ne peut pas supporter $r+1$ pertes consécutives, la probabilité de ruine sur une séquence est :

\[
q^{r+1}.
\]

Dans le cas équilibré :

\[
q=\frac12.
\]

Donc :

\[
\Prob(\text{ruine})=\left(\frac12\right)^{r+1}.
\]

\subsection{Espérance avec capital limité}

Le joueur :
\begin{itemize}
    \item gagne $1$ avec probabilité $1-q^{r+1}$ ;
    \item perd $2^{r+1}-1$ avec probabilité $q^{r+1}$.
\end{itemize}

L'espérance est donc :

\[
\E=1(1-q^{r+1})-(2^{r+1}-1)q^{r+1}.
\]

Dans le cas équilibré :

\[
q=\frac12.
\]

Donc :

\[
\E=1\left(1-\frac{1}{2^{r+1}}\right)
-
(2^{r+1}-1)\frac{1}{2^{r+1}}.
\]

Or :

\[
(2^{r+1}-1)\frac{1}{2^{r+1}}
=
1-\frac{1}{2^{r+1}}.
\]

Donc :

\[
\E=0.
\]

La martingale ne crée pas d'espérance positive dans un jeu équilibré.

Dans un jeu défavorable, l'espérance devient négative.

\subsection{Conclusion}

\begin{conclusionbox}{martingale}
La martingale ne bat pas le hasard.

Elle transforme un risque fréquent et modéré en un risque rare mais catastrophique.

En pratique, elle échoue à cause :
\begin{itemize}
    \item du capital limité du joueur ;
    \item des plafonds de mise ;
    \item des longues séries de pertes ;
    \item de l'espérance négative des jeux de casino.
\end{itemize}
\end{conclusionbox}

\begin{oral}{rédaction}
La martingale donne l'impression d'être infaillible car une victoire finit par récupérer toutes les pertes précédentes. Mais les mises suivent une suite géométrique : après $r$ pertes, le joueur a déjà perdu $2^r-1$. Le capital nécessaire augmente donc extrêmement vite. Avec un capital limité, une longue série de pertes suffit à ruiner le joueur. La martingale ne supprime pas le risque, elle le concentre dans un événement rare mais très coûteux.
\end{oral}

%=========================================================
\section{Piste 6 -- Jouer tous les jours jusqu'à l'infini : le robot joueur}

\subsection{Contexte concret}

\begin{contexte}{robot joueur}
On imagine un robot qui joue automatiquement chaque jour à un jeu d'argent.

Il n'a pas d'émotions, pas d'intuition, pas de superstition. Il applique simplement la même stratégie tous les jours.

La question est :

\[
\text{Si le robot joue jusqu'à l'infini, finit-il forcément gagnant ?}
\]
\end{contexte}

\subsection{Idée générale}

\begin{idee}{loi des grands nombres}
À très long terme, le gain moyen par partie se rapproche de l'espérance du jeu.
\end{idee}

\subsection{Modélisation}

On note :

\[
X_i
\]

le gain du robot le jour $i$.

On suppose que les variables $X_i$ sont :
\begin{itemize}
    \item indépendantes ;
    \item de même loi ;
    \item d'espérance $\mu$.
\end{itemize}

Le gain total après $n$ jours est :

\[
S_n=X_1+\cdots+X_n.
\]

Le gain moyen par jour est :

\[
M_n=\frac{S_n}{n}.
\]

\subsection{Mise en équation}

Par linéarité de l'espérance :

\[
\E(S_n)=n\mu.
\]

Donc :

\[
\E(M_n)=\E\left(\frac{S_n}{n}\right).
\]

Par linéarité :

\[
\E(M_n)=\frac{1}{n}\E(S_n).
\]

Donc :

\[
\E(M_n)=\mu.
\]

\subsection{Loi des grands nombres}

La loi des grands nombres dit que :

\[
M_n=\frac{S_n}{n}
\]

se rapproche de $\mu$ lorsque $n$ devient très grand.

On écrit, dans une version avancée :

\[
\frac{S_n}{n}\xrightarrow[n\to+\infty]{\Prob}\mu.
\]

Cela signifie que le gain moyen observé se rapproche de l'espérance théorique.

\subsection{Résolution dans un jeu défavorable}

Si :

\[
\mu<0,
\]

alors :

\[
\frac{S_n}{n}\approx \mu
\]

pour $n$ grand.

Donc :

\[
S_n\approx n\mu.
\]

Comme :

\[
n\mu\to -\infty,
\]

le robot perd de plus en plus d'argent au total.

\subsection{Conclusion}

\begin{conclusionbox}{robot infini}
Un robot qui joue indéfiniment à un jeu d'espérance négative ne finit pas par battre le hasard.

Le long terme ne corrige pas l'injustice du jeu : il la révèle.

Plus le robot joue, plus son gain moyen se rapproche de l'espérance négative.
\end{conclusionbox}

\begin{oral}{rédaction}
On peut modéliser le robot par une suite de variables aléatoires indépendantes $X_i$, chacune représentant le gain d'une journée. Le gain total est $S_n=X_1+\cdots+X_n$, et le gain moyen est $\frac{S_n}{n}$. La loi des grands nombres affirme que cette moyenne se rapproche de l'espérance $\mu$. Donc si le jeu est défavorable, c'est-à-dire si $\mu<0$, le robot perd en moyenne de plus en plus d'argent au fil du temps.
\end{oral}

%=========================================================
\section{Piste 7 -- Continuité des probabilités : finir par gagner au moins une fois}

\subsection{Contexte concret}

\begin{contexte}{``si je joue assez longtemps, je vais bien gagner''}
Un joueur dit : ``Même si la probabilité de gagner est faible, si je joue assez longtemps, je finirai bien par gagner.''

Cette phrase contient une part de vérité, mais elle est souvent mal interprétée.
\end{contexte}

\subsection{Idée générale}

\begin{idee}{au moins un succès}
Si la probabilité de gagner une partie est positive, alors la probabilité de gagner au moins une fois parmi $n$ parties tend vers $1$ quand $n$ tend vers l'infini.
\end{idee}

\subsection{Modélisation}

On suppose que la probabilité de gagner une partie est :

\[
p>0.
\]

Les parties sont indépendantes.

On note $A_n$ l'événement :

\[
A_n=\text{``obtenir au moins une victoire parmi les }n\text{ premières parties''}.
\]

\subsection{Mise en équation}

L'événement contraire est :

\[
\overline{A_n}=\text{``perdre les }n\text{ premières parties''}.
\]

La probabilité de perdre une partie est :

\[
1-p.
\]

Comme les parties sont indépendantes :

\[
\Prob(\overline{A_n})=(1-p)^n.
\]

Donc :

\[
\Prob(A_n)=1-(1-p)^n.
\]

\subsection{Résolution}

Comme :

\[
0<1-p<1,
\]

alors :

\[
(1-p)^n\to 0.
\]

Donc :

\[
\Prob(A_n)\to 1.
\]

\subsection{Interprétation rigoureuse}

Cela signifie :

\[
\lim_{n\to+\infty}\Prob(A_n)=1.
\]

Autrement dit, si on joue indéfiniment, la probabilité de gagner au moins une fois tend vers $1$.

Mais cela ne dit rien sur le bilan financier total.

\subsection{Conclusion}

\begin{conclusionbox}{victoire ponctuelle contre rentabilité}
Il est vrai qu'en jouant très longtemps, on finit probablement par gagner au moins une fois.

Mais cette victoire ponctuelle peut coûter beaucoup plus cher que ce qu'elle rapporte.

Il faut distinguer :
\begin{itemize}
    \item la probabilité de gagner au moins une fois ;
    \item l'espérance du gain total.
\end{itemize}
\end{conclusionbox}

\begin{oral}{rédaction}
Si chaque partie a une probabilité de succès $p>0$, alors la probabilité de ne jamais gagner pendant les $n$ premières parties est $(1-p)^n$. Donc la probabilité de gagner au moins une fois est $1-(1-p)^n$, qui tend vers $1$. Cela explique pourquoi l'affirmation ``je finirai bien par gagner'' est partiellement vraie. Mais ce résultat ne garantit pas que le joueur sera rentable : il peut gagner une fois après avoir déjà perdu beaucoup plus.
\end{oral}

\begin{prepa}{continuité croissante}
Les événements $A_n$ forment une suite croissante :

\[
A_1\subset A_2\subset A_3\subset\cdots.
\]

La continuité croissante des probabilités affirme que :

\[
\Prob\left(\bigcup_{n\geq1}A_n\right)
=
\lim_{n\to+\infty}\Prob(A_n).
\]

Ici, l'événement :

\[
\bigcup_{n\geq1}A_n
\]

signifie : ``gagner au moins une fois au cours d'un nombre fini de parties''.
\end{prepa}

%=========================================================
\section{Piste 8 -- Blackjack : optimiser une décision avec l'espérance conditionnelle}

\subsection{Contexte concret}

\begin{contexte}{décision au blackjack}
Au blackjack, le joueur ne se contente pas d'attendre le hasard. Il doit prendre des décisions : tirer une carte, s'arrêter, doubler, séparer certaines paires.

Contrairement à la roulette, il existe donc une part d'optimisation mathématique.
\end{contexte}

\subsection{Idée générale}

\begin{idee}{comparer deux espérances conditionnelles}
Dans une situation donnée, la décision optimale est celle dont l'espérance conditionnelle est la plus grande.
\end{idee}

\subsection{Modélisation simplifiée}

On suppose que le joueur a un total :

\[
s.
\]

Il peut :
\begin{itemize}
    \item s'arrêter ;
    \item tirer une carte.
\end{itemize}

On note $C$ la valeur de la prochaine carte.

Si :

\[
s+C>21,
\]

le joueur perd immédiatement.

\subsection{Probabilité de dépasser 21}

Le joueur dépasse 21 si :

\[
C>21-s.
\]

Donc :

\[
\Prob(\text{dépasser})=\Prob(C>21-s).
\]

\subsection{Exemple avec un total de 16}

Si :

\[
s=16,
\]

alors le joueur dépasse si :

\[
16+C>21.
\]

Donc :

\[
C>5.
\]

Dans un modèle simplifié où les cartes $1,2,\dots,10$ sont équiprobables :

\[
\Prob(C>5)=\frac{5}{10}=\frac12.
\]

\subsection{Comparaison des décisions}

On note :

\[
G_T
\]

le gain si le joueur tire, et :

\[
G_S
\]

le gain s'il s'arrête.

La décision mathématique consiste à comparer :

\[
\E(G_T\mid s)
\]

et :

\[
\E(G_S\mid s).
\]

Si :

\[
\E(G_T\mid s)>\E(G_S\mid s),
\]

alors il vaut mieux tirer.

Si :

\[
\E(G_T\mid s)<\E(G_S\mid s),
\]

alors il vaut mieux s'arrêter.

\subsection{Conclusion}

\begin{conclusionbox}{blackjack}
Le blackjack est un jeu où certaines décisions peuvent être optimisées.

Mais optimiser une décision ne signifie pas forcément rendre le jeu favorable.

La stratégie optimale réduit l'avantage du casino, mais ne le supprime pas toujours.
\end{conclusionbox}

\begin{oral}{rédaction}
Au blackjack, on peut utiliser l'espérance conditionnelle. Par exemple, si le joueur a $16$, il dépasse $21$ dès que la prochaine carte vaut au moins $6$. Dans un modèle simplifié, cela représente une probabilité de $\frac12$. La stratégie consiste alors à comparer l'espérance du gain si l'on tire avec l'espérance du gain si l'on s'arrête. Le blackjack montre que les mathématiques peuvent améliorer une décision, sans garantir un gain à long terme.
\end{oral}

\begin{prepa}{programmation dynamique}
On peut définir une fonction $V(s)$ représentant l'espérance maximale du joueur lorsqu'il possède un total $s$.

On peut écrire :

\[
V(s)=\max\left(V_{\text{stop}}(s),V_{\text{tirer}}(s)\right).
\]

Cette relation est une idée de programmation dynamique : une décision optimale se construit en comparant les valeurs futures possibles.
\end{prepa}

%=========================================================
\section{Piste 9 -- Ruine du joueur : pourquoi le capital compte autant que la probabilité}

\subsection{Contexte concret}

\begin{contexte}{joueur contre casino}
Même si un jeu était parfaitement équilibré, le joueur a souvent un capital limité, alors que le casino possède un capital très important.

Cette asymétrie change tout.
\end{contexte}

\subsection{Idée générale}

\begin{idee}{marche aléatoire}
Le capital du joueur peut être modélisé comme une marche aléatoire : il monte ou descend d'une unité à chaque partie.
\end{idee}

\subsection{Modélisation}

Le joueur commence avec un capital :

\[
a.
\]

Il s'arrête s'il atteint :

\[
N
\]

ou s'il tombe à :

\[
0.
\]

À chaque partie :
\[
\begin{cases}
+1 & \text{avec probabilité }p,\\
-1 & \text{avec probabilité }q=1-p.
\end{cases}
\]

Dans le cas équilibré :

\[
p=q=\frac12.
\]

\subsection{Résultat dans le cas équilibré}

Dans le cas équilibré, la probabilité d'atteindre $N$ avant $0$ en partant de $a$ est :

\[
\Prob(\text{atteindre }N)=\frac{a}{N}.
\]

Donc :

\[
\Prob(\text{ruine})=1-\frac{a}{N}.
\]

\subsection{Justification intuitive}

Si le jeu est équilibré, le capital moyen ne doit pas avoir tendance à monter ou descendre.

On peut alors montrer que la probabilité de succès est proportionnelle à la position initiale $a$ entre les deux barrières $0$ et $N$.

Si $a$ est proche de $0$, la ruine est très probable.

Si $a$ est proche de $N$, atteindre $N$ est plus probable.

\subsection{Conclusion}

\begin{conclusionbox}{capital fini}
Même dans un jeu équilibré, un joueur avec un petit capital face à un objectif élevé a une forte probabilité de ruine.

Si :

\[
N\gg a,
\]

alors :

\[
\frac{a}{N}\approx 0,
\]

donc :

\[
\Prob(\text{ruine})\approx 1.
\]

Le casino n'a donc pas besoin que chaque partie soit extrêmement défavorable : son avantage vient aussi du grand nombre de parties et de la différence de capital.
\end{conclusionbox}

\begin{oral}{rédaction}
Le capital du joueur peut être vu comme une marche aléatoire. Dans le cas équilibré, s'il part de $a$ et vise $N$, la probabilité d'atteindre $N$ avant la ruine est $\frac{a}{N}$. Si le casino a un capital beaucoup plus grand que celui du joueur, cette probabilité est très faible. Cela montre que le risque de ruine dépend non seulement des probabilités du jeu, mais aussi du capital disponible.
\end{oral}

%=========================================================
\section{Piste 10 -- Jeux avec gains continus : pourquoi les intégrales apparaissent}

\subsection{Contexte concret}

\begin{contexte}{gain variable}
Dans certains jeux ou paris, le gain n'est pas seulement une liste de quelques valeurs. On peut modéliser un gain comme une quantité continue.

Par exemple, un bonus peut dépendre d'un coefficient aléatoire entre $0$ et $100$.
\end{contexte}

\subsection{Idée générale}

\begin{idee}{espérance continue}
Quand une variable peut prendre toutes les valeurs d'un intervalle, on remplace les sommes par des intégrales.
\end{idee}

\subsection{Définition}

\begin{definitionbox}{densité de probabilité}
Une variable aléatoire continue $X$ peut être décrite par une densité $f$.

La probabilité que $X$ appartienne à l'intervalle $[a,b]$ vaut :

\[
\Prob(a\leq X\leq b)=\int_a^b f(x)\,dx.
\]

La densité vérifie :

\[
f(x)\geq 0
\]

et :

\[
\int_{-\infty}^{+\infty}f(x)\,dx=1.
\]
\end{definitionbox}

\subsection{Espérance continue}

\begin{definitionbox}{espérance d'une variable continue}
Si $X$ admet une densité $f$, son espérance est :

\[
\E(X)=\int_{-\infty}^{+\infty}x f(x)\,dx.
\]
\end{definitionbox}

\subsection{Modélisation}

On suppose que le gain brut $X$ est uniforme sur $[0,100]$.

Cela signifie que :

\[
f(x)=\frac{1}{100}
\]

pour :

\[
0\leq x\leq 100.
\]

\subsection{Résolution}

\[
\E(X)=\int_0^{100}x\frac{1}{100}\,dx.
\]

Donc :

\[
\E(X)=\frac{1}{100}\left[\frac{x^2}{2}\right]_0^{100}.
\]

Ainsi :

\[
\E(X)=\frac{1}{100}\cdot \frac{10000}{2}=50.
\]

Si le ticket coûte $60$ euros, le gain algébrique est :

\[
Y=X-60.
\]

Donc :

\[
\E(Y)=\E(X)-60=50-60=-10.
\]

\subsection{Conclusion}

\begin{conclusionbox}{intégrale}
L'intégrale généralise l'idée de moyenne pondérée.

Même lorsque le gain est continu, l'espérance permet de juger la rentabilité du jeu.

Ici, le joueur perd en moyenne $10$ euros.
\end{conclusionbox}

\begin{oral}{rédaction}
Lorsque les gains possibles ne sont pas discrets, on utilise une densité de probabilité. L'espérance se calcule alors avec une intégrale. Dans l'exemple d'un gain uniforme entre $0$ et $100$ euros, le gain moyen est $50$ euros. Si le ticket coûte $60$ euros, l'espérance du gain net est $-10$ euros. Le jeu reste donc défavorable, même si le gain maximal paraît attractif.
\end{oral}

%=========================================================
\section{Piste 11 -- Critère de Kelly : optimiser sa mise quand le jeu est favorable}

\subsection{Contexte concret}

\begin{contexte}{pari légèrement favorable}
Supposons qu'un joueur ait trouvé un pari réellement favorable. Par exemple, il gagne avec une probabilité supérieure à $\frac12$ et le gain est égal à la mise.

Doit-il miser tout son capital ?

Intuitivement non, car une perte totale l'empêcherait de continuer.
\end{contexte}

\subsection{Idée générale}

\begin{idee}{croissance du capital}
Le critère de Kelly cherche à maximiser la croissance à long terme du capital, pas seulement l'espérance du gain d'une partie.
\end{idee}

\subsection{Modélisation}

Le joueur possède un capital $C$.

Il mise une fraction $f$ de ce capital.

Avec probabilité $p$, il gagne sa mise et son capital devient :

\[
C(1+f).
\]

Avec probabilité $q=1-p$, il perd sa mise et son capital devient :

\[
C(1-f).
\]

\subsection{Pourquoi le logarithme ?}

Sur plusieurs parties, le capital est multiplié par des facteurs successifs.

Si les facteurs sont :

\[
A_1,A_2,\dots,A_n,
\]

alors :

\[
C_n=C_0A_1A_2\cdots A_n.
\]

En prenant le logarithme :

\[
\ln(C_n)=\ln(C_0)+\ln(A_1)+\cdots+\ln(A_n).
\]

Le logarithme transforme un produit en somme, ce qui permet d'utiliser une moyenne.

\subsection{Mise en équation}

On maximise :

\[
G(f)=p\ln(1+f)+q\ln(1-f).
\]

\subsection{Résolution}

On dérive :

\[
G'(f)=\frac{p}{1+f}-\frac{q}{1-f}.
\]

On cherche :

\[
G'(f)=0.
\]

Donc :

\[
\frac{p}{1+f}=\frac{q}{1-f}.
\]

On croise :

\[
p(1-f)=q(1+f).
\]

Donc :

\[
p-pf=q+qf.
\]

Ainsi :

\[
p-q=f(p+q).
\]

Comme :

\[
p+q=1,
\]

on obtient :

\[
f=p-q.
\]

Or :

\[
q=1-p.
\]

Donc :

\[
f=2p-1.
\]

\subsection{Conclusion}

\begin{conclusionbox}{Kelly}
Si :

\[
p\leq\frac12,
\]

alors :

\[
f\leq 0.
\]

La meilleure décision est de ne pas jouer.

Si :

\[
p>\frac12,
\]

alors le critère donne une fraction positive du capital à miser.

Même dans un jeu favorable, il ne faut pas miser tout son capital.
\end{conclusionbox}

\begin{oral}{rédaction}
Le critère de Kelly répond à une question différente : si un pari est favorable, quelle fraction du capital faut-il miser ? Comme le capital se multiplie au fil des parties, on maximise l'espérance du logarithme du capital. Dans le modèle simple où l'on gagne ou perd sa mise, le calcul donne $f=2p-1$. Si le jeu n'est pas favorable, c'est-à-dire si $p\leq\frac12$, alors la fraction optimale est nulle : il vaut mieux ne pas jouer.
\end{oral}

%=========================================================
\section{Piste 12 -- Illusion du joueur : après plusieurs pertes, la victoire est-elle plus probable ?}

\subsection{Contexte concret}

\begin{contexte}{roulette}
À la roulette, si le rouge n'est pas sorti depuis longtemps, certains joueurs pensent qu'il est ``dû'' de sortir.

C'est une erreur classique appelée illusion du joueur.
\end{contexte}

\subsection{Idée générale}

\begin{idee}{indépendance des parties}
Si les parties sont indépendantes, les résultats passés ne changent pas la probabilité du prochain résultat.
\end{idee}

\subsection{Modélisation}

On note :

\[
A=\text{``rouge au prochain lancer''}.
\]

On note :

\[
B=\text{``les cinq derniers lancers n'ont pas donné rouge''}.
\]

Si les lancers sont indépendants :

\[
\Prob(A\mid B)=\Prob(A).
\]

\subsection{Définition de la probabilité conditionnelle}

\begin{definitionbox}{probabilité conditionnelle}
Si $\Prob(B)>0$, la probabilité de $A$ sachant $B$ est :

\[
\Prob(A\mid B)=\frac{\Prob(A\cap B)}{\Prob(B)}.
\]

Si $A$ et $B$ sont indépendants, alors :

\[
\Prob(A\cap B)=\Prob(A)\Prob(B).
\]

Donc :

\[
\Prob(A\mid B)=\frac{\Prob(A)\Prob(B)}{\Prob(B)}=\Prob(A).
\]
\end{definitionbox}

\subsection{Conclusion}

\begin{conclusionbox}{illusion du joueur}
Après une série de pertes, la probabilité de gagner ne devient pas automatiquement plus grande.

Dans un jeu indépendant, le passé ne corrige pas le futur.

La loi des grands nombres ne signifie pas que le hasard compense immédiatement les déséquilibres passés.
\end{conclusionbox}

\begin{oral}{rédaction}
L'illusion du joueur consiste à croire qu'après plusieurs pertes, une victoire devient plus probable. Mais si les parties sont indépendantes, la probabilité conditionnelle du prochain succès reste égale à la probabilité initiale. Mathématiquement, $\Prob(A\mid B)=\Prob(A)$. Le hasard ne cherche pas à équilibrer les résultats à court terme.
\end{oral}
%=========================================================
\section{Piste 13 -- La martingale à la roulette : technique gagnante ou illusion mathématique ?}

\subsection{Contexte concret}

\begin{contexte}{un joueur face à la roulette}
Un joueur arrive au casino avec une idée simple : il mise toujours sur rouge. S'il gagne, il empoche son gain. S'il perd, il double sa mise au coup suivant. Son raisonnement est le suivant :

\begin{quote}
    \og Si je finis par gagner une fois, je récupère toutes mes pertes précédentes et je gagne une unité. \fg{}
\end{quote}

Cette technique est appelée \textbf{martingale classique}. Elle est très connue dans les jeux d'argent, notamment à la roulette, car elle donne l'impression de transformer un jeu de hasard en méthode presque sûre.

La question est donc :

\[
\text{La martingale permet-elle réellement de gagner à la roulette ?}
\]
\end{contexte}

\subsection{Définition de la technique}

\begin{definitionbox}{martingale de casino}
Dans le langage des jeux, une martingale est une stratégie de mise qui consiste à modifier sa mise en fonction des résultats précédents.

La martingale classique consiste à :
\begin{itemize}
    \item miser une unité au premier coup ;
    \item si on perd, doubler la mise ;
    \item continuer à doubler après chaque perte ;
    \item dès qu'on gagne, revenir à la mise initiale.
\end{itemize}

Ainsi, les mises successives après des pertes sont :

\[
1,\;2,\;4,\;8,\;16,\;32,\dots
\]

Autrement dit, la mise après \(r\) pertes consécutives vaut :

\[
2^r.
\]
\end{definitionbox}

\begin{prepa}{lien avec la notion mathématique de martingale}
En probabilités avancées, une martingale est une suite de variables aléatoires dont l'espérance conditionnelle future, sachant le passé, est égale à la valeur présente. Autrement dit, dans un jeu équitable, connaître les gains passés ne donne pas d'avantage pour le gain futur. Une définition classique en temps discret est :
\[
\mathbb{E}(M_{n+1}\mid M_0,M_1,\dots,M_n)=M_n.
\]
Cette idée vient historiquement du vocabulaire des jeux : une martingale désignait déjà une stratégie de mise. Le texte fourni rappelle aussi cette origine et donne la définition probabiliste :
\[
\mathbb{E}(X_{t+1}\mid X_0,\dots,X_t)=X_t.
\]
:contentReference[oaicite:1]{index=1}
\end{prepa}

\subsection{Pourquoi la martingale semble fonctionner ?}

\begin{idee}{récupération des pertes}
La martingale semble fonctionner parce qu'après une série de pertes, une seule victoire permet de récupérer toutes les pertes précédentes et de gagner une unité.
\end{idee}

Supposons que le joueur perde \(r\) fois de suite.

Il a alors perdu :

\[
1+2+4+\cdots+2^{r-1}.
\]

C'est une somme géométrique. On sait que :

\[
1+2+4+\cdots+2^{r-1}=2^r-1.
\]

Au coup suivant, il mise :

\[
2^r.
\]

S'il gagne ce coup, il reçoit un gain net égal à sa mise, car le pari rouge/noir paie \(1\) contre \(1\). Son bilan total est alors :

\[
\text{gain du dernier coup} - \text{pertes précédentes}
=
2^r-(2^r-1).
\]

Donc :

\[
2^r-(2^r-1)=1.
\]

\begin{equationbox}{bilan d'un cycle gagnant}
Quel que soit le nombre de pertes précédentes, si le joueur finit par gagner avant d'être bloqué, son gain net est :

\[
\boxed{1}
\]
\end{equationbox}

\begin{oral}{rédaction possible}
La martingale paraît convaincante car les pertes forment une somme géométrique. Après \(r\) pertes, le joueur a perdu \(2^r-1\). Il mise alors \(2^r\). S'il gagne, son gain du dernier coup compense exactement toutes les pertes précédentes et laisse un bénéfice de \(1\). C'est cette propriété qui donne l'illusion d'une méthode sûre.
\end{oral}

\subsection{Mais à la roulette, le pari rouge/noir n'est pas équitable}

\begin{contexte}{le rôle du zéro}
À la roulette, miser sur rouge ne signifie pas avoir une chance sur deux de gagner.

En roulette européenne, il y a :
\[
18 \text{ cases rouges},\quad 18 \text{ cases noires},\quad 1 \text{ case verte }0.
\]

Il y a donc \(37\) cases au total.

En roulette américaine, il y a :
\[
18 \text{ cases rouges},\quad 18 \text{ cases noires},\quad 2 \text{ cases vertes }0\text{ et }00.
\]

Il y a donc \(38\) cases au total.
\end{contexte}

\subsubsection*{Cas 1 : roulette européenne}

Si le joueur mise \(1\) euro sur rouge, alors :

\[
X=
\begin{cases}
1 & \text{avec probabilité } \frac{18}{37},\\
-1 & \text{avec probabilité } \frac{19}{37}.
\end{cases}
\]

En effet, le joueur gagne sur les \(18\) cases rouges, mais perd sur les \(18\) cases noires et sur le zéro.

L'espérance est :

\[
\mathbb{E}(X)
=
1\times \frac{18}{37}
+
(-1)\times \frac{19}{37}.
\]

Donc :

\[
\mathbb{E}(X)
=
\frac{18}{37}-\frac{19}{37}
=
-\frac{1}{37}.
\]

Ainsi :

\[
\mathbb{E}(X)\approx -0{,}027.
\]

Le joueur perd donc en moyenne environ \(2{,}70\%\) de sa mise à chaque coup en roulette européenne.

\begin{equationbox}{avantage casino en roulette européenne}
Sur une mise de \(1\), l'espérance du joueur est :

\[
\boxed{\mathbb{E}(X)=-\frac{1}{37}\approx -2{,}70\%}
\]
\end{equationbox}

\subsubsection*{Cas 2 : roulette américaine avec \(0\) et \(00\)}

À la roulette américaine, il y a \(38\) cases :

\[
18 \text{ rouges},\quad 18 \text{ noires},\quad 0,\quad 00.
\]

Donc, sur un pari rouge :

\[
X=
\begin{cases}
1 & \text{avec probabilité } \frac{18}{38},\\
-1 & \text{avec probabilité } \frac{20}{38}.
\end{cases}
\]

L'espérance vaut :

\[
\mathbb{E}(X)
=
1\times \frac{18}{38}
+
(-1)\times \frac{20}{38}.
\]

Donc :

\[
\mathbb{E}(X)
=
\frac{18}{38}-\frac{20}{38}
=
-\frac{2}{38}.
\]

Ainsi :

\[
\mathbb{E}(X)=-\frac{1}{19}\approx -0{,}0526.
\]

Le joueur perd donc en moyenne environ \(5{,}26\%\) de sa mise à chaque coup.

\begin{equationbox}{avantage casino en roulette américaine}
Sur une mise de \(1\), l'espérance du joueur est :

\[
\boxed{\mathbb{E}(X)=-\frac{1}{19}\approx -5{,}26\%}
\]
\end{equationbox}

\begin{oral}{rédaction possible}
La martingale repose souvent sur l'idée que rouge/noir est un pari à \(50\%\). Mais à la roulette, ce n'est pas vrai. En roulette européenne, il y a 18 cases rouges, 18 cases noires et un zéro. Donc la probabilité de gagner sur rouge est seulement \(\frac{18}{37}\), et non \(\frac12\). L'espérance d'une mise de \(1\) euro vaut alors \(-\frac{1}{37}\). En roulette américaine, avec le double zéro, la probabilité de gagner tombe à \(\frac{18}{38}\), et l'espérance devient \(-\frac{1}{19}\). Le double zéro augmente donc fortement l'avantage du casino.
\end{oral}

\subsection{Ce que le casino a fait pour rendre la martingale inutile}

\subsubsection*{1. Le zéro et le double zéro rendent le jeu défavorable}

\begin{idee}{l'avantage structurel du casino}
Le zéro, puis le double zéro dans la roulette américaine, empêchent le pari rouge/noir d'être équitable.
\end{idee}

Si le pari rouge/noir était parfaitement équitable, on aurait :

\[
p=\frac12,
\qquad
q=\frac12.
\]

Mais en roulette européenne :

\[
p=\frac{18}{37}<\frac12,
\qquad
q=\frac{19}{37}>\frac12.
\]

En roulette américaine :

\[
p=\frac{18}{38}<\frac12,
\qquad
q=\frac{20}{38}>\frac12.
\]

Le joueur a donc plus de chances de perdre que de gagner à chaque coup.

\begin{conclusionbox}{rôle du zéro}
Le zéro et le double zéro ne bloquent pas directement la martingale, mais ils rendent chaque coup défavorable.

La martingale ne change pas cette espérance négative : elle ne fait que modifier la manière dont les pertes apparaissent.
\end{conclusionbox}

\subsubsection*{2. Les plafonds de mise bloquent la progression exponentielle}

\begin{idee}{la vraie arme contre la martingale}
La vraie protection du casino contre la martingale est le plafond de mise : le joueur n'a pas le droit de doubler indéfiniment.
\end{idee}

Supposons que la mise minimale soit \(1\) euro et que la mise maximale autorisée soit \(M\).

Le joueur peut miser :

\[
1,\;2,\;4,\;8,\dots,2^r
\]

tant que :

\[
2^r\leq M.
\]

Donc le nombre maximal de pertes successives qu'il peut suivre est donné par :

\[
r\leq \log_2(M).
\]

Par exemple, si :

\[
M=512,
\]

alors :

\[
512=2^9.
\]

Le joueur peut donc aller jusqu'à une mise de \(512\), mais ne peut pas miser \(1024\) au coup suivant.

Or, pour encaisser une série de \(10\) pertes consécutives, il faudrait pouvoir miser :

\[
1024.
\]

Le plafond de mise casse donc la martingale.

\begin{equationbox}{plafond de mise}
Si la mise maximale est \(M\), la martingale est bloquée dès que :

\[
2^r > M.
\]

Autrement dit, le nombre de doubles possibles est limité par :

\[
\boxed{r\leq \log_2(M)}
\]
\end{equationbox}

\subsubsection*{3. Le capital du joueur est fini}

Même sans plafond de mise, le joueur n'a pas un capital infini.

Après \(r\) pertes consécutives, la somme totale déjà perdue est :

\[
1+2+4+\cdots+2^{r-1}=2^r-1.
\]

Pour continuer après \(r\) pertes, il faut en plus pouvoir miser :

\[
2^r.
\]

Il faut donc disposer au total de :

\[
(2^r-1)+2^r=2^{r+1}-1.
\]

\begin{equationbox}{capital nécessaire}
Pour survivre à \(r\) pertes puis tenter de se refaire au coup suivant, il faut un capital d'au moins :

\[
\boxed{2^{r+1}-1}
\]
\end{equationbox}

Par exemple :

\[
\begin{array}{c|c|c}
\text{Nombre de pertes }r & \text{mise suivante} & \text{capital nécessaire} \\
\hline
5 & 32 & 63 \\
10 & 1024 & 2047 \\
15 & 32768 & 65535 \\
20 & 1048576 & 2097151
\end{array}
\]

La croissance est exponentielle.

\begin{conclusionbox}{capital fini}
Le joueur peut supporter quelques pertes, mais pas une longue série.

La martingale transforme donc de nombreuses petites victoires probables en une perte rare mais énorme.
\end{conclusionbox}

\subsection{Probabilité d'être bloqué par une série de pertes}

Supposons que le joueur puisse supporter \(r\) pertes consécutives, mais pas \(r+1\).

La probabilité de perdre \(r+1\) fois de suite est :

\[
q^{r+1},
\]

où \(q\) est la probabilité de perdre un coup.

\subsubsection*{À la roulette européenne}

On a :

\[
q=\frac{19}{37}.
\]

Donc :

\[
\Prob(\text{série fatale})
=
\left(\frac{19}{37}\right)^{r+1}.
\]

\subsubsection*{À la roulette américaine}

On a :

\[
q=\frac{20}{38}=\frac{10}{19}.
\]

Donc :

\[
\Prob(\text{série fatale})
=
\left(\frac{10}{19}\right)^{r+1}.
\]

Comme :

\[
\frac{10}{19}>\frac{19}{37},
\]

la série fatale est plus probable en roulette américaine.

\begin{resolution}{exemple numérique}
Si le joueur peut supporter \(9\) pertes, mais pas \(10\), alors la série fatale est une série de \(10\) pertes consécutives.

En roulette européenne :

\[
\Prob(\text{10 pertes})=
\left(\frac{19}{37}\right)^{10}
\approx 0{,}00127.
\]

Cela fait environ :

\[
0{,}127\%.
\]

En roulette américaine :

\[
\Prob(\text{10 pertes})=
\left(\frac{10}{19}\right)^{10}
\approx 0{,}00163.
\]

Cela fait environ :

\[
0{,}163\%.
\]

Ces probabilités semblent faibles, mais la perte associée est très grande.
\end{resolution}

\subsection{Espérance d'un cycle de martingale avec capital limité}

On suppose que le joueur peut supporter \(r\) pertes, mais pas \(r+1\).

Il y a deux issues :
\begin{itemize}
    \item il gagne avant la série fatale et gagne \(1\) unité ;
    \item il perd \(r+1\) fois de suite et perd toutes les mises engagées.
\end{itemize}

La perte dans le cas fatal est :

\[
1+2+4+\cdots+2^r=2^{r+1}-1.
\]

La probabilité du cas fatal est :

\[
q^{r+1}.
\]

La probabilité de gagner le cycle est donc :

\[
1-q^{r+1}.
\]

L'espérance du cycle est :

\[
\mathbb{E}
=
1\cdot(1-q^{r+1})
-
(2^{r+1}-1)q^{r+1}.
\]

On simplifie :

\[
\mathbb{E}
=
1-q^{r+1}-(2^{r+1}-1)q^{r+1}.
\]

Donc :

\[
\mathbb{E}
=
1-2^{r+1}q^{r+1}.
\]

Ainsi :

\[
\boxed{\mathbb{E}=1-(2q)^{r+1}}
\]

\subsubsection*{Cas d'un jeu parfaitement équitable}

Si le jeu était équitable :

\[
q=\frac12.
\]

Alors :

\[
2q=1.
\]

Donc :

\[
\mathbb{E}=1-1^{r+1}=0.
\]

Même dans un jeu parfaitement équitable, la martingale ne crée pas de gain moyen.

\subsubsection*{Cas de la roulette européenne}

En roulette européenne :

\[
q=\frac{19}{37}.
\]

Donc :

\[
2q=\frac{38}{37}>1.
\]

Ainsi :

\[
\mathbb{E}
=
1-\left(\frac{38}{37}\right)^{r+1}<0.
\]

La martingale est perdante en espérance.

\subsubsection*{Cas de la roulette américaine}

En roulette américaine :

\[
q=\frac{20}{38}=\frac{10}{19}.
\]

Donc :

\[
2q=\frac{20}{19}>1.
\]

Ainsi :

\[
\mathbb{E}
=
1-\left(\frac{20}{19}\right)^{r+1}<0.
\]

La martingale est encore plus défavorable.

\begin{conclusionbox}{résultat central}
La martingale ne change pas l'espérance du jeu.

Si le jeu est équitable, elle reste d'espérance nulle.

Si le jeu est défavorable, elle reste défavorable.

À la roulette, comme le zéro ou le double zéro donnent un avantage au casino, la martingale est perdante en moyenne.
\end{conclusionbox}

\subsection{Pourquoi le double zéro rend la technique encore moins intéressante}

\begin{idee}{double zéro}
Le double zéro augmente la probabilité de perdre un pari rouge/noir.
\end{idee}

À la roulette européenne :

\[
q_E=\frac{19}{37}\approx 0{,}5135.
\]

À la roulette américaine :

\[
q_A=\frac{20}{38}\approx 0{,}5263.
\]

Donc :

\[
q_A>q_E.
\]

La probabilité d'une série de pertes est donc plus grande en roulette américaine :

\[
q_A^{r+1}>q_E^{r+1}.
\]

De plus, l'espérance d'un cycle devient plus négative :

\[
1-(2q_A)^{r+1}
<
1-(2q_E)^{r+1}.
\]

\begin{conclusionbox}{effet du double zéro}
Le double zéro n'est pas seulement une case en plus.

Il augmente :
\begin{itemize}
    \item la probabilité de perdre un coup ;
    \item la probabilité d'une longue série de pertes ;
    \item l'avantage du casino ;
    \item la perte moyenne d'une martingale avec capital limité.
\end{itemize}

Il rend donc la martingale encore plus inefficace.
\end{conclusionbox}

\subsection{Synthèse sous forme de tableau}

\begin{center}
\renewcommand{\arraystretch}{1.5}
\begin{tabular}{p{4cm}p{4cm}p{4cm}}
\toprule
\textbf{Situation} & \textbf{Roulette européenne} & \textbf{Roulette américaine} \\
\midrule
Nombre de cases & \(37\) & \(38\) \\
Cases vertes & \(0\) & \(0\) et \(00\) \\
Probabilité de gagner rouge & \(\frac{18}{37}\) & \(\frac{18}{38}\) \\
Probabilité de perdre rouge & \(\frac{19}{37}\) & \(\frac{20}{38}\) \\
Espérance d'une mise de \(1\) & \(-\frac{1}{37}\) & \(-\frac{1}{19}\) \\
Avantage casino & \(2{,}70\%\) & \(5{,}26\%\) \\
Effet sur la martingale & Défavorable & Encore plus défavorable \\
\bottomrule
\end{tabular}
\end{center}

\subsection{Conclusion générale de la partie}

\begin{conclusionbox}{réponse à la question}
La martingale à la roulette est une stratégie séduisante mais mathématiquement perdante en pratique.

Elle semble garantir un petit gain, mais seulement si trois conditions irréalistes sont réunies :
\begin{itemize}
    \item le joueur possède un capital infini ;
    \item le casino n'impose aucun plafond de mise ;
    \item le jeu est au moins équitable.
\end{itemize}

Or, dans un vrai casino :
\begin{itemize}
    \item le capital du joueur est fini ;
    \item les tables ont des mises maximales ;
    \item la roulette contient un zéro, voire un double zéro ;
    \item le pari rouge/noir est donc défavorable.
\end{itemize}

Le casino n'a donc pas besoin de prévoir chaque stratégie individuellement : la combinaison du zéro, du double zéro, des plafonds de mise et du capital fini suffit à rendre la martingale inutile.
\end{conclusionbox}

\begin{oral}{rédaction finale possible}
La martingale repose sur une idée simple : doubler après chaque perte pour qu'une seule victoire récupère toutes les pertes précédentes. Mathématiquement, cela fonctionne seulement si le joueur peut doubler indéfiniment. Or les mises nécessaires augmentent comme des puissances de \(2\), ce qui impose rapidement un capital énorme. En plus, les casinos imposent des plafonds de mise, ce qui empêche de continuer la stratégie après une longue série de pertes. Enfin, à la roulette, rouge/noir n'est pas un pari équitable à cause du zéro, et encore moins en roulette américaine avec le double zéro. L'espérance reste donc négative. La martingale ne supprime pas le risque : elle le transforme en une perte rare mais très importante.
\end{oral}
%=========================================================
%=========================================================
\section{Comparaison de deux joueurs : mise fixe contre martingale}

\subsection{Contexte concret}

\begin{contexte}{deux stratégies face à la roulette}
On compare deux joueurs qui jouent à la roulette sur les paris simples rouge/noir.

\begin{itemize}
    \item Le joueur A mise toujours la même somme \(U_0\), et choisit aléatoirement rouge ou noir à chaque tour.
    \item Le joueur B utilise une martingale : il commence par miser \(U_0\), puis double sa mise après chaque perte. Dès qu'il gagne, il revient à la mise initiale \(U_0\).
\end{itemize}

On cherche à répondre à la question :

\[
\text{Quelle stratégie est la plus rentable après plusieurs tours ?}
\]

Pour traiter le problème proprement, on modélise les gains et pertes mathématiquement.
\end{contexte}

%---------------------------------------------------------
\subsection{Probabilité de gain sur rouge/noir}

On note \(p\) la probabilité de gagner un coup sur un pari rouge/noir, et \(q\) la probabilité de perdre.

Ainsi :

\[
q=1-p.
\]

À la roulette européenne :

\[
p=\frac{18}{37},
\qquad
q=\frac{19}{37}.
\]

À la roulette américaine :

\[
p=\frac{18}{38},
\qquad
q=\frac{20}{38}.
\]

Dans les deux cas :

\[
p<q.
\]

Donc le jeu est défavorable au joueur.

\begin{attention}{choisir rouge ou noir au hasard ne change rien}
Si le joueur choisit rouge ou noir au hasard, cela ne change pas sa probabilité de gagner.

En effet, rouge et noir ont le même nombre de cases. Le problème ne vient pas du choix rouge/noir, mais de la présence du zéro, ou du double zéro, qui fait perdre le joueur sur un pari rouge/noir.
\end{attention}

%=========================================================
\subsection{Joueur A : mise fixe \(U_0\)}

\subsubsection*{Modélisation}

Le joueur A mise toujours :

\[
U_0
\]

à chaque tour.

On note \(X_k\) le gain algébrique du joueur A au \(k\)-ième tour.

Alors :

\[
X_k=
\begin{cases}
+U_0 & \text{avec probabilité }p,\\
-U_0 & \text{avec probabilité }q.
\end{cases}
\]

Le gain total après \(n\) tours est :

\[
S_n=X_1+X_2+\cdots+X_n.
\]

\subsubsection*{Espérance d'un tour}

On calcule :

\[
\mathbb{E}(X_k)
=
U_0p+(-U_0)q.
\]

Donc :

\[
\mathbb{E}(X_k)
=
U_0(p-q).
\]

Comme \(q=1-p\), on peut aussi écrire :

\[
\mathbb{E}(X_k)=U_0(2p-1).
\]

À la roulette européenne :

\[
\mathbb{E}(X_k)
=
U_0\left(\frac{18}{37}-\frac{19}{37}\right)
=
-\frac{U_0}{37}.
\]

À la roulette américaine :

\[
\mathbb{E}(X_k)
=
U_0\left(\frac{18}{38}-\frac{20}{38}\right)
=
-\frac{2U_0}{38}
=
-\frac{U_0}{19}.
\]

\begin{equationbox}{espérance d'un coup à mise fixe}
Pour une roulette européenne :

\[
\boxed{\mathbb{E}(X_k)=-\frac{U_0}{37}}
\]

Pour une roulette américaine :

\[
\boxed{\mathbb{E}(X_k)=-\frac{U_0}{19}}
\]
\end{equationbox}

\subsubsection*{Espérance après \(n\) tours}

Par linéarité de l'espérance :

\[
\mathbb{E}(S_n)
=
\mathbb{E}(X_1+\cdots+X_n).
\]

Donc :

\[
\mathbb{E}(S_n)
=
\mathbb{E}(X_1)+\cdots+\mathbb{E}(X_n).
\]

Comme les variables \(X_k\) ont toutes la même espérance :

\[
\mathbb{E}(S_n)=n\mathbb{E}(X_1).
\]

Ainsi :

\[
\mathbb{E}(S_n)=nU_0(p-q).
\]

À la roulette européenne :

\[
\boxed{\mathbb{E}(S_n)=-\frac{nU_0}{37}}
\]

À la roulette américaine :

\[
\boxed{\mathbb{E}(S_n)=-\frac{nU_0}{19}}
\]

\subsubsection*{Expression du gain selon le nombre de victoires}

Si le joueur gagne \(W_n\) fois parmi les \(n\) tours, alors il perd \(n-W_n\) fois.

Son gain total est :

\[
S_n=U_0W_n-U_0(n-W_n).
\]

Donc :

\[
S_n=U_0(2W_n-n).
\]

Comme \(W_n\) suit une loi binomiale :

\[
W_n\sim \mathcal{B}(n,p),
\]

on a :

\[
\mathbb{E}(W_n)=np.
\]

Donc :

\[
\mathbb{E}(S_n)
=
U_0(2np-n)
=
nU_0(2p-1),
\]

ce qui correspond bien au résultat précédent.

\begin{conclusionbox}{joueur A}
Le joueur A perd en moyenne une fraction constante de chaque mise.

Sa perte moyenne est proportionnelle au nombre de tours :

\[
\mathbb{E}(S_n)=nU_0(p-q)<0.
\]

Il n'y a pas d'effet explosif : les mises restent constantes, donc le risque est contrôlé, mais la stratégie reste perdante en moyenne.
\end{conclusionbox}

%=========================================================
\subsection{Joueur B : martingale avec mise initiale \(U_0\)}

\subsubsection*{Principe de la martingale}

Le joueur B commence par miser :

\[
U_0.
\]

S'il perd, il mise :

\[
2U_0.
\]

S'il perd encore, il mise :

\[
4U_0.
\]

Après \(r\) pertes consécutives, il mise :

\[
2^rU_0.
\]

Dès qu'il gagne, il revient à la mise initiale \(U_0\).

\subsubsection*{Gain ou perte après une série de pertes}

Supposons que le joueur perde \(r\) fois de suite.

Il a alors perdu :

\[
U_0+2U_0+4U_0+\cdots+2^{r-1}U_0.
\]

On factorise :

\[
U_0(1+2+4+\cdots+2^{r-1}).
\]

Or :

\[
1+2+\cdots+2^{r-1}=2^r-1.
\]

Donc la perte après \(r\) pertes consécutives est :

\[
U_0(2^r-1).
\]

Au coup suivant, le joueur mise :

\[
2^rU_0.
\]

S'il gagne, son gain sur ce coup est :

\[
2^rU_0.
\]

Son bilan total du cycle est donc :

\[
2^rU_0-U_0(2^r-1).
\]

Donc :

\[
2^rU_0-U_0(2^r-1)=U_0.
\]

\begin{equationbox}{gain d'un cycle de martingale réussi}
Si le joueur finit par gagner avant d'être bloqué, alors son gain net est toujours :

\[
\boxed{U_0}
\]
\end{equationbox}

\subsubsection*{Perte après \(n\) pertes consécutives}

Si le joueur perd les \(n\) premiers tours de la martingale, alors il a misé :

\[
U_0,\;2U_0,\;4U_0,\;\dots,\;2^{n-1}U_0.
\]

Sa perte totale vaut :

\[
L_n=U_0(1+2+4+\cdots+2^{n-1}).
\]

Donc :

\[
\boxed{L_n=U_0(2^n-1)}
\]

Cette formule est essentielle : la perte croît exponentiellement avec \(n\).

Par exemple :

\[
\begin{array}{c|c|c}
n & \text{mise au }n\text{-ième tour} & \text{perte totale si }n\text{ pertes} \\
\hline
1 & U_0 & U_0 \\
2 & 2U_0 & 3U_0 \\
3 & 4U_0 & 7U_0 \\
4 & 8U_0 & 15U_0 \\
5 & 16U_0 & 31U_0 \\
10 & 512U_0 & 1023U_0 \\
\end{array}
\]

\subsubsection*{Expression du gain selon le rang de la première victoire}

On note \(T\) le rang de la première victoire.

Par exemple :
\[
T=1
\]
signifie que le joueur gagne tout de suite.

\[
T=3
\]
signifie qu'il perd les deux premiers coups puis gagne au troisième.

Si la première victoire arrive au rang \(k\), alors le joueur a perdu \(k-1\) fois puis gagné au \(k\)-ième coup.

Son gain net vaut :

\[
G_k=U_0.
\]

Donc :

\[
\boxed{T=k \Longrightarrow G_k=U_0}
\]

tant que le joueur peut continuer à doubler.

\subsubsection*{Probabilité de la première victoire au rang \(k\)}

Pour que la première victoire ait lieu au rang \(k\), il faut :
\begin{itemize}
    \item perdre les \(k-1\) premiers coups ;
    \item gagner le \(k\)-ième.
\end{itemize}

La probabilité est donc :

\[
\mathbb{P}(T=k)=q^{k-1}p.
\]

\subsubsection*{Cas réaliste : capital fini ou plafond de mise}

Dans la réalité, le joueur ne peut pas doubler indéfiniment.

On suppose qu'il peut jouer au maximum \(n\) coups dans un cycle de martingale.

Cela signifie :
\begin{itemize}
    \item s'il gagne parmi les \(n\) premiers coups, il gagne \(U_0\) ;
    \item s'il perd les \(n\) coups, il subit une perte de \(U_0(2^n-1)\).
\end{itemize}

La probabilité de perdre les \(n\) coups est :

\[
q^n.
\]

La probabilité de gagner au moins une fois parmi les \(n\) coups est :

\[
1-q^n.
\]

Donc le gain \(M_n\) du cycle de martingale limité à \(n\) coups vaut :

\[
M_n=
\begin{cases}
U_0 & \text{avec probabilité }1-q^n,\\
-U_0(2^n-1) & \text{avec probabilité }q^n.
\end{cases}
\]

\subsubsection*{Espérance du cycle de martingale limité}

On calcule :

\[
\mathbb{E}(M_n)
=
U_0(1-q^n)-U_0(2^n-1)q^n.
\]

On factorise par \(U_0\) :

\[
\mathbb{E}(M_n)
=
U_0\left(1-q^n-(2^n-1)q^n\right).
\]

On simplifie :

\[
1-q^n-(2^n-1)q^n
=
1-q^n-2^nq^n+q^n.
\]

Donc :

\[
\mathbb{E}(M_n)
=
U_0(1-2^nq^n).
\]

Ainsi :

\[
\boxed{\mathbb{E}(M_n)=U_0\left(1-(2q)^n\right)}
\]

\subsubsection*{Interprétation}

Si le jeu était parfaitement équitable :

\[
q=\frac12.
\]

Alors :

\[
2q=1.
\]

Donc :

\[
\mathbb{E}(M_n)=U_0(1-1^n)=0.
\]

Même dans un jeu équitable, la martingale ne crée pas de rentabilité moyenne.

À la roulette européenne :

\[
q=\frac{19}{37}.
\]

Donc :

\[
2q=\frac{38}{37}>1.
\]

Ainsi :

\[
\mathbb{E}(M_n)
=
U_0\left(1-\left(\frac{38}{37}\right)^n\right)<0.
\]

À la roulette américaine :

\[
q=\frac{20}{38}=\frac{10}{19}.
\]

Donc :

\[
2q=\frac{20}{19}>1.
\]

Ainsi :

\[
\mathbb{E}(M_n)
=
U_0\left(1-\left(\frac{20}{19}\right)^n\right)<0.
\]

\begin{conclusionbox}{joueur B}
La martingale donne une forte probabilité de gagner \(U_0\), mais elle expose à une faible probabilité de perdre une somme énorme :

\[
U_0(2^n-1).
\]

Son espérance reste négative à la roulette, car \(q>\frac12\).

La martingale n'améliore donc pas la rentabilité moyenne : elle modifie seulement la répartition du risque.
\end{conclusionbox}

%=========================================================
\subsection{Comparaison des deux méthodes}

\subsubsection*{Joueur A : mise fixe}

Après \(n\) tours, son gain espéré est :

\[
\mathbb{E}(S_n)=nU_0(p-q).
\]

À la roulette européenne :

\[
\mathbb{E}(S_n)=-\frac{nU_0}{37}.
\]

À la roulette américaine :

\[
\mathbb{E}(S_n)=-\frac{nU_0}{19}.
\]

\subsubsection*{Joueur B : martingale limitée à \(n\) coups}

Sur un cycle limité à \(n\) coups, son gain espéré est :

\[
\mathbb{E}(M_n)=U_0\left(1-(2q)^n\right).
\]

À la roulette européenne :

\[
\mathbb{E}(M_n)
=
U_0\left(1-\left(\frac{38}{37}\right)^n\right).
\]

À la roulette américaine :

\[
\mathbb{E}(M_n)
=
U_0\left(1-\left(\frac{20}{19}\right)^n\right).
\]

\subsubsection*{Attention : les deux comparaisons ne portent pas exactement sur la même mise totale}

Pour le joueur A, après \(n\) tours, la somme totale misée est :

\[
nU_0.
\]

Pour le joueur B, s'il perd \(n\) fois, la somme totale engagée peut atteindre :

\[
U_0(2^n-1).
\]

Donc la martingale peut engager beaucoup plus d'argent que la mise fixe.

Pour comparer honnêtement, il faut donc regarder :
\begin{itemize}
    \item l'espérance de gain ;
    \item le risque maximal ;
    \item la somme totale potentiellement engagée ;
    \item la probabilité de ruine.
\end{itemize}

%---------------------------------------------------------
\subsection{Quelle méthode est la plus rentable ?}

\begin{conclusionbox}{réponse mathématique}
Aucune des deux méthodes n'est rentable à la roulette.

\[
\boxed{\text{Les deux stratégies ont une espérance négative.}}
\]

La mise fixe est plus simple et limite les pertes, car le joueur ne mise jamais plus que \(U_0\) par tour.

La martingale donne l'impression de gagner souvent, mais elle expose à une perte rare et énorme.

Donc :
\[
\boxed{\text{la martingale n'est pas plus rentable ; elle est seulement plus risquée.}}
\]
\end{conclusionbox}

\subsection{Interprétation en termes de risque}

\begin{center}
\renewcommand{\arraystretch}{1.5}
\begin{tabular}{p{4cm}p{5cm}p{5cm}}
\toprule
\textbf{Stratégie} & \textbf{Avantage apparent} & \textbf{Problème mathématique} \\
\midrule
Mise fixe \(U_0\) & Risque contrôlé, pertes progressives & Espérance négative à cause du zéro \\
\midrule
Martingale & Beaucoup de cycles gagnants de \(U_0\) & Perte rare mais exponentielle : \(U_0(2^n-1)\) \\
\bottomrule
\end{tabular}
\end{center}

\subsection{Rédaction possible à l'oral}

\begin{oral}{comparaison des deux joueurs}
On peut comparer deux stratégies. Le premier joueur mise toujours la même somme \(U_0\). Son gain après \(n\) tours est une somme de variables aléatoires valant \(+U_0\) ou \(-U_0\). Son espérance est donc \(nU_0(p-q)\), qui est négative à la roulette car la probabilité de perdre est légèrement supérieure à la probabilité de gagner.

Le second joueur utilise une martingale. S'il gagne avant d'être bloqué, il gagne toujours \(U_0\). Mais s'il perd \(n\) fois de suite, sa perte vaut \(U_0(2^n-1)\), ce qui augmente exponentiellement. L'espérance d'un cycle limité à \(n\) coups vaut \(U_0(1-(2q)^n)\), qui est négative dès que \(q>\frac12\). La martingale n'est donc pas plus rentable que la mise fixe : elle transforme des pertes régulières en un risque rare mais très violent.
\end{oral}

\end{document}
```
